\section{Round-three positive closure: an analytic--convex commutation theorem}
\label{sec:r3-d1}

This section gives D1 an independent theorem.  It proves when the microscopic
large-deviation limit commutes with differentiation, source-dependent
conditioning, convex duality, continuous contraction, Gaussian likelihood
ratios, and nonlinear semigroup tangents.  The B-series supplies a concrete
hard-sphere instance of its hypotheses.

\subsection{The normalized finite-volume hierarchy}

Let \(X\) be a separable weighted source Banach space and let
\(\mathcal X_\varepsilon\) be an \(X^*\)-valued path observable under
\(\mathbb P_\varepsilon\).  Let
\(\mu_\varepsilon\to\infty\), and define
\[
 Q_\varepsilon(\Theta)
 =\frac1{\mu_\varepsilon}
 \log\mathbb E_\varepsilon
 e^{\mu_\varepsilon\langle\Theta,
 \mathcal X_\varepsilon\rangle}.
\]
For a finite-dimensional constraint map
\(C:X^*\to\mathbb R^k\), write \(C^*\lambda\in X\).

\begin{proposition}[Exact speed normalization]
\label{prop:r3-d1-normalization}
For \(F,G\in X\),
\[
 DQ_\varepsilon(\Theta)[F]
 =\mathbb E_{\Theta,\varepsilon}
 \langle F,\mathcal X_\varepsilon\rangle,
\]
\[
 D^2Q_\varepsilon(\Theta)[F,G]
 =\mu_\varepsilon
 \operatorname{Cov}_{\Theta,\varepsilon}
 \left(
 \langle F,\mathcal X_\varepsilon\rangle,
 \langle G,\mathcal X_\varepsilon\rangle
 \right).
\]
Thus the limiting Hessian is the covariance of the
\(\sqrt{\mu_\varepsilon}\)-fluctuation field, not the ordinary covariance of
an unscaled limiting macroscopic variable.
\end{proposition}

\begin{proof}
Differentiate the finite logarithmic moment generating function once and
twice.  The second derivative of the exponent contributes one factor
\(\mu_\varepsilon\), while the outer normalization removes only one of the two
factors.
\end{proof}

\subsection{The commutation hypotheses}

Fix a real source ball \(B\Subset X\) and a complex neighborhood
\(B_{\mathbb C}^+\).  The hypotheses are:
\begin{enumerate}[label=(H\arabic*)]
\item \(Q_\varepsilon\) extends holomorphically to \(B_{\mathbb C}^+\), is
      locally uniformly bounded there, and converges pointwise on a set with
      an accumulation point to \(Q\);
\item the laws of \(\mathcal X_\varepsilon\) are exponentially tight in a
      weighted weak topology and have a good exposed-point lower bound;
\item on the constraint directions,
\[
 cI_k\le C D^2Q(\Theta)C^*\le CI_k;
\]
\item every observable contraction used below has a continuous adjoint on
      sources, or is the limit of exponentially good continuous
      approximations.
\end{enumerate}
In the hard-sphere application, (H1) is B2, (H2) is the B2 joint LDP, (H3) is
B1/B3 strict constraint covariance, and (H4) is the weighted weak typing of
the chosen density--collision observables.

\begin{theorem}[Analytic--convex commutation theorem]
\label{thm:r3-d1-commutation}
Under (H1)--(H4), the following statements hold on every smaller real source
ball.
\begin{enumerate}[label=(\roman*)]
\item \(Q_\varepsilon\to Q\) locally uniformly with all Fr\'echet derivatives.
\item The good rate is the convex dual
\[
 I(x)=\sup_{\Theta\in X}
 \{\langle\Theta,x\rangle-Q(\Theta)+Q(0)\}
\]
on the exposed regular domain, with lower-semicontinuous extension elsewhere.
\item For target \(a\), the constrained pressure is
\[
 Q^a(\Theta)
 =\inf_\lambda
 \{Q(\Theta+C^*\lambda)-\lambda\cdot a\}
 -\inf_\lambda
 \{Q(C^*\lambda)-\lambda\cdot a\},
\]
and the infimum has a unique analytic source-dependent minimizer
\(\lambda_\Theta\).
\item Its Hessian is the Schur complement
\[
 D^2Q^a
 =Q_{\Theta\Theta}
 -Q_{\Theta\lambda}Q_{\lambda\lambda}^{-1}
 Q_{\lambda\Theta}.
\]
\item If \(T:X^*\to E\) is a typed continuous contraction, then
\[
 Q_T(\eta)=Q(T^*\eta),
 \qquad
 I_T(y)=\inf_{Tx=y}I(x),
\]
and first and second derivatives commute with \(T\).
\item The constrained rates and contracted rates commute:
\[
 \inf_{Tx=y,\ Cx=a}I(x)
 =\left((Q^a\circ T^*)^*\right)(y)
\]
on the regular exposed domain, followed by lower-semicontinuous closure.
\end{enumerate}
\end{theorem}

\begin{proof}
Local boundedness and pointwise convergence of holomorphic functions give
normal-family convergence.  Uniqueness of the real limit identifies every
subsequence, and Cauchy's formula gives local uniform convergence of all
derivatives, proving (i).

Exponential tightness gives the pressure upper bound.  Exponential tilting at
an exposed source, together with the exposed-point lower bound, gives the
matching lower bound.  Approximation by exposed points and lower-semicontinuous
closure yield (ii).

For (iii), strict constraint covariance makes
\(\lambda\mapsto Q(\Theta+C^*\lambda)-\lambda\cdot a\) uniformly strictly
convex.  Its critical equation is
\[
 C DQ(\Theta+C^*\lambda_\Theta)=a.
\]
The analytic inverse-function theorem gives the unique analytic minimizer.
Differentiating this equation and then the envelope gives (iv).

For (v),
\[
 \langle\eta,T\mathcal X_\varepsilon\rangle
 =\langle T^*\eta,\mathcal X_\varepsilon\rangle,
\]
so the pressure identity is exact at finite \(\varepsilon\).  Duality gives
the contraction rate, and the chain rule gives its derivatives.  If \(T\) is
only exponentially approximable, the extended contraction theorem and normal
source convergence give the same result.  Finally, minimizing \(I\) under the
joint affine map \((T,C)\) and taking its source dual yields (vi); strict
constraint convexity permits the \(C\)-dual to be eliminated first.
\end{proof}

\subsection{Central limits and canonical likelihood ratios}

Fix a regular source \(\Theta\), let
\[
 m_\Theta=DQ(\Theta),
 \qquad
 \Sigma_\Theta=D^2Q(\Theta),
\]
and put
\[
 Z_\varepsilon
 =\sqrt{\mu_\varepsilon}
 (\mathcal X_\varepsilon-m_\Theta).
\]

\begin{theorem}[Gaussian tangent and likelihood-ratio convergence]
\label{thm:r3-d1-likelihood}
Under the hypotheses of
Theorem~\ref{thm:r3-d1-commutation}, every finite-dimensional projection of
\(Z_\varepsilon\) converges to a centered Gaussian field \(Z\) with covariance
\(\Sigma_\Theta\).  For \(h\) in such a finite-dimensional source space, the
local canonical likelihood ratio
\[
 L_\varepsilon(h)
 =\exp\left\{
 \sqrt{\mu_\varepsilon}\langle h,
 \mathcal X_\varepsilon-m_\Theta\rangle
 -\mu_\varepsilon\left[
 Q_\varepsilon(\Theta+h/\sqrt{\mu_\varepsilon})
 -Q_\varepsilon(\Theta)
 -\frac1{\sqrt{\mu_\varepsilon}}DQ_\varepsilon(\Theta)h
 \right]
 \right\}
\]
converges jointly with \(Z_\varepsilon\) to
\[
 L(h)=\exp\left\{\langle h,Z\rangle
 -\tfrac12\Sigma_\Theta(h,h)\right\}.
\]
The likelihood ratios are uniformly integrable.  Under \(L_\varepsilon(h)\),
the Gaussian limit is shifted by \(\Sigma_\Theta h\).
\end{theorem}

\begin{proof}
Expand the analytic pressure at
\(\Theta+it/\sqrt{\mu_\varepsilon}\).  The quadratic term converges to
\(-\frac12\Sigma_\Theta(t,t)\), while Cauchy bounds make the remainder
\(O(\mu_\varepsilon^{-1/2}|t|^3)\).  L\'evy's theorem gives the Gaussian
limit.  The same expansion at real \(h\), jointly with imaginary test sources,
gives the displayed likelihood limit.  Analyticity on a larger real ball gives
an \(L^{1+\delta}\) bound, hence uniform integrability.  Multiplying the joint
characteristic function by \(L(h)\) gives the Gaussian mean shift.
\end{proof}

\begin{corollary}[Conditioned Gaussian tangent]
\label{cor:r3-d1-conditioned}
After source-dependent microcanonical conditioning, the Gaussian covariance is
\[
 \Sigma^a
 =\Sigma-\Sigma C^*(C\Sigma C^*)^{-1}C\Sigma.
\]
Every typed projection \(T\) has covariance
\(T\Sigma^aT^*\), and its local likelihood ratios are the corresponding
Cameron--Martin exponentials.
\end{corollary}

\begin{proof}
Apply part (iv) of Theorem~\ref{thm:r3-d1-commutation} and then
Theorem~\ref{thm:r3-d1-likelihood} to the constrained pressure.  Apply the
chain rule for \(T\).
\end{proof}

\subsection{Dynamic action, Hessian, and semigroup in one diagram}

Let \(Q_{s,t}\) be an additive path pressure and let \(I_{s,t}\) be its B2
convex dual action.  Let \(S_{s,t}\) be the B4 Lax--Oleinik semigroup.

\begin{theorem}[Rate--pressure--semigroup triangle]
\label{thm:r3-d1-triangle}
On the regular short-time hard-sphere domain, the following operations
commute:
\[
 \begin{array}{ccc}
 Q_{s,t} & \xrightarrow{\text{convex dual}} & I_{s,t}\\
 \downarrow D,D^2 && \downarrow\text{quadratic tangent}\\
 (m,\Sigma) & \xrightarrow{\text{Legendre}} &
 \tfrac12\|\cdot\|_{\Sigma^{-1}}^2,
 \end{array}
\]
and
\[
 (S_{s,t}\Phi)(f)
 =\sup_{\rho_s=f}
 \{\Phi(\rho_t)-I_{s,t}(\rho)\}.
\]
The dynamic additivity of \(I\) is equivalent to the nonlinear tower of
\(S\).  After a finite-dimensional Gaussian contraction, the likelihood
ratios of Theorem~\ref{thm:r3-d1-likelihood} converge to the Girsanov
Cameron--Martin density, and the exponential transform of the Gaussian
semigroup has the quadratic theta generator.
\end{theorem}

\begin{proof}
The first row is B2/B3 duality.  Analytic second-order expansion and
Theorem~\ref{thm:r3-d1-likelihood} show that the local Legendre dual converges
to the inverse-covariance quadratic form on the covariance range.  B4
constructs the variational semigroup from the additive action; path
concatenation proves both additivity and the tower.  Under a finite-dimensional
projection, the limiting local likelihood is the Cameron--Martin exponential.
Applying the logarithmic exponential transform to the Gaussian generator adds
\(\frac\theta2D\Phi^TAD\Phi\).
\end{proof}

\subsection{Typed contraction registry}

Every entry below includes its parent platform and map.
\begin{enumerate}[label=(\alph*)]
\item \textbf{HSBG density--collision:} the identity map on the B2 weighted
      pair; pressure, action, and Gaussian tangent are B2--B4.
\item \textbf{HSBG finite observables:} any weighted continuous linear map of
      density and actual collision flow; rates and covariances follow from
      Theorem~\ref{thm:r3-d1-commutation}.
\item \textbf{HSBG prepared rays:} finite-dimensional constraints are treated
      by the source-dependent saddle and the Schur complement, never by a
      fixed zero-source phase.
\item \textbf{Tagged equilibrium packet:} a tagged-particle Brownian theorem
      belongs to its separately scaled equilibrium hard-sphere family; once
      its convergence map is supplied, its Gaussian pressure and likelihood
      diagrams are an instance of the contraction theorem, not a finite-time
      identity with the present kinetic path.
\item \textbf{Heat-bath and Lorentz packets:} these retain their augmented or
      fixed-scatterer platform labels.  They enter the coproduct only through
      their own proved microscopic laws and scaling maps.
\end{enumerate}
No item is obtained by deleting particles from a different microscopic law.

\subsection{Calibrated scalar rays}

\begin{theorem}[Mechanical ray calibration]
\label{thm:r3-d1-calibration}
Let a scalar constrained observable \(G\) have differentiable rate
\(I_G\), and set \(\theta=I_G'(a)\).  The canonical preparation law uses the
likelihood exponent \(\theta G\).  A smooth, cash-additive, strongly
time-consistent certainty equivalent is entropic with some constant
\(\vartheta\).  If its likelihood cocycle is required to be this same
canonical preparation cocycle, then \(\vartheta=\theta\).  Without the
cocycle requirement, \(\vartheta\) is an external parameter.
\end{theorem}

\begin{proof}
The scalar convex saddle gives the canonical exponent \(\theta\).  Entropic
rigidity gives the exponential form with constant \(\vartheta\).  Equality of
the two Radon--Nikodym cocycles on the nonconstant generator \(G\) forces
\(\vartheta=\theta\).
\end{proof}

\begin{theorem}[Round-three D1 closure]
\label{thm:r3-d1-main}
For the hard-sphere source domain proved in B1--B4, the
Boltzmann--Grad limit commutes with pressure differentiation,
source-dependent microcanonical conditioning, convex duality, typed
contraction, Gaussian tangent formation, and canonical likelihood-ratio
limits.  The exact second derivative is the speed-normalized covariance, and
all downstream Gaussian, Girsanov, and theta diagrams use that normalization.
\end{theorem}

\begin{proof}
Apply Theorems~\ref{thm:r3-d1-commutation},
\ref{thm:r3-d1-likelihood}, \ref{thm:r3-d1-triangle}, and
\ref{thm:r3-d1-calibration} to the B2 analytic pressure, the B1 saddle, the B3
joint covariance, and the B4 action semigroup.
\end{proof}
